\[ %%%%%%%%%%%%%%%%%%%%%%%%%%%%%%%%%%%%%%%%%%%%%%%%%%%%%%%%%%%%%%%%%%%%%%%%%%%%%%%% \subsection{Adaptive Stability Control Module} \label{sec:stability_control} %%%%%%%%%%%%%%%%%%%%%%%%%%%%%%%%%%%%%%%%%%%%%%%%%%%%%%%%%%%%%%%%%%%%%%%%%%%%%%%% Long-horizon prediction of nonlinear dynamical systems is fundamentally limited by error accumulation during recursive forecasting. Even small deviations introduced by a learned evolution operator may propagate through successive prediction steps, resulting in trajectory divergence and loss of physical consistency. To address this challenge, PHYSGEN introduces an Adaptive Stability Control Module (ASCM), designed to monitor, regulate, and correct the evolution of latent physical states during long-term prediction. The central hypothesis of ASCM is that stability should not be treated only as a consequence of model training, but as an active computational process integrated into the prediction architecture. %%%%%%%%%%%%%%%%%%%%%%%%%%%%%%%%%%%%%%%%%%%%%%%%%%%%%%%%%%%%%%%%%%%%%%%%%%%%%%%% \subsubsection{Stability-Aware State Evolution} %%%%%%%%%%%%%%%%%%%%%%%%%%%%%%%%%%%%%%%%%%%%%%%%%%%%%%%%%%%%%%%%%%%%%%%%%%%%%%%% The standard latent evolution process can be written as: \begin{equation} H_{t+1} = \Psi_{\theta}(H_t,G_t,P_t). \end{equation} In PHYSGEN, this evolution is extended with an adaptive stability correction: \begin{equation} H_{t+1} = \Psi_{\theta}(H_t,G_t,P_t) + C_{stab}(H_t,G_t), \end{equation} where $C_{stab}$ represents the stability control operator. The stability controller evaluates the current state trajectory and determines whether corrective regulation is required. %%%%%%%%%%%%%%%%%%%%%%%%%%%%%%%%%%%%%%%%%%%%%%%%%%%%%%%%%%%%%%%%%%%%%%%%%%%%%%%% \subsubsection{Stability State Estimation} %%%%%%%%%%%%%%%%%%%%%%%%%%%%%%%%%%%%%%%%%%%%%%%%%%%%%%%%%%%%%%%%%%%%%%%%%%%%%%%% The stability controller receives information from the latent trajectory: \begin{equation} s_t = S_{\theta} ( H_t, H_{t-1}, \Delta H_t ), \end{equation} where: \begin{itemize} \item $s_t$ represents the estimated stability state; \item $\Delta H_t$ describes the recent latent state variation; \item $S_\theta$ denotes the stability estimation network. \end{itemize} The stability state may capture: \begin{itemize} \item growth of prediction error; \item abnormal latent-space deviation; \item violation of physical invariants; \item excessive energy growth; \item instability of temporal evolution. \end{itemize} %%%%%%%%%%%%%%%%%%%%%%%%%%%%%%%%%%%%%%%%%%%%%%%%%%%%%%%%%%%%%%%%%%%%%%%%%%%%%%%% \subsubsection{Adaptive Correction Mechanism} %%%%%%%%%%%%%%%%%%%%%%%%%%%%%%%%%%%%%%%%%%%%%%%%%%%%%%%%%%%%%%%%%%%%%%%%%%%%%%%% Based on the estimated stability state, PHYSGEN applies an adaptive correction: \begin{equation} C_{stab} = -\lambda_s \nabla_H \mathcal{L}_{stab}, \end{equation} where $\lambda_s$ represents the stability regulation coefficient. The stability loss is defined as: \begin{equation} \mathcal{L}_{stab} = \| H_{t+1}-H_t \|^2 + \beta \mathcal{C}_{phys}(H_t), \end{equation} where: \begin{itemize} \item the first term controls excessive latent state changes; \item $\mathcal{C}_{phys}$ measures physical inconsistency. \end{itemize} Therefore, the controller prevents unrealistic deviations while preserving necessary nonlinear dynamics. %%%%%%%%%%%%%%%%%%%%%%%%%%%%%%%%%%%%%%%%%%%%%%%%%%%%%%%%%%%%%%%%%%%%%%%%%%%%%%%% \subsubsection{Lyapunov-Inspired Stability Regulation} %%%%%%%%%%%%%%%%%%%%%%%%%%%%%%%%%%%%%%%%%%%%%%%%%%%%%%%%%%%%%%%%%%%%%%%%%%%%%%%% The design of ASCM is inspired by stability analysis of nonlinear dynamical systems. A Lyapunov function $V(H)$ describes the stability of a system trajectory: \begin{equation} \Delta V = V(H_{t+1})-V(H_t). \end{equation} For stable evolution, the desired condition is: \begin{equation} \Delta V \leq 0. \end{equation} PHYSGEN does not require an explicitly known analytical Lyapunov function. Instead, the stability controller learns a latent stability representation: \begin{equation} V_\theta(H_t), \end{equation} and regulates the trajectory toward dynamically stable regions: \begin{equation} V_\theta(H_{t+1}) \leq V_\theta(H_t)+\epsilon. \end{equation} This formulation provides a flexible approximation of stability principles for complex systems where analytical analysis is infeasible. %%%%%%%%%%%%%%%%%%%%%%%%%%%%%%%%%%%%%%%%%%%%%%%%%%%%%%%%%%%%%%%%%%%%%%%%%%%%%%%% \subsubsection{Adaptive Time Evolution Control} %%%%%%%%%%%%%%%%%%%%%%%%%%%%%%%%%%%%%%%%%%%%%%%%%%%%%%%%%%%%%%%%%%%%%%%%%%%%%%%% Different physical regimes may require different prediction behaviors. For example, smooth evolution and highly nonlinear transitions should not be treated identically. Therefore, PHYSGEN introduces an adaptive evolution coefficient: \begin{equation} \eta_t = g_\theta(s_t), \end{equation} which controls the magnitude of latent updates: \begin{equation} H_{t+1} = H_t + \eta_t \Delta H_t. \end{equation} The coefficient $\eta_t$ allows the model to: \begin{itemize} \item accelerate evolution in stable regimes; \item reduce updates near unstable states; \item improve robustness during chaotic transitions. \end{itemize} %%%%%%%%%%%%%%%%%%%%%%%%%%%%%%%%%%%%%%%%%%%%%%%%%%%%%%%%%%%%%%%%%%%%%%%%%%%%%%%% \subsubsection{Long-Horizon Rollout with Stability Control} %%%%%%%%%%%%%%%%%%%%%%%%%%%%%%%%%%%%%%%%%%%%%%%%%%%%%%%%%%%%%%%%%%%%%%%%%%%%%%%% During multi-step prediction, the complete PHYSGEN rollout becomes: \begin{equation} H_{t+k} = \prod_{i=0}^{k-1} \left( \Psi_\theta + C_{stab} \right) H_t . \end{equation} The objective is not only minimizing short-term prediction error: \begin{equation} \mathcal{L}_{pred} = \| X_t-\hat{X}_t \|, \end{equation} but maintaining physically consistent trajectories: \begin{equation} \mathcal{L}_{total} = \mathcal{L}_{pred} + \lambda_s\mathcal{L}_{stab} + \lambda_p\mathcal{L}_{phys}. \end{equation} %%%%%%%%%%%%%%%%%%%%%%%%%%%%%%%%%%%%%%%%%%%%%%%%%%%%%%%%%%%%%%%%%%%%%%%%%%%%%%%% \subsubsection{Summary} %%%%%%%%%%%%%%%%%%%%%%%%%%%%%%%%%%%%%%%%%%%%%%%%%%%%%%%%%%%%%%%%%%%%%%%%%%%%%%%% The Adaptive Stability Control Module transforms long-horizon forecasting from a passive prediction problem into an actively regulated dynamical process. By continuously monitoring latent evolution, estimating stability conditions, and applying adaptive corrections, PHYSGEN reduces error accumulation and improves the reliability of recursive prediction. This mechanism represents a key distinction between PHYSGEN and conventional graph-based forecasting models, where stability is typically considered an emergent property of training rather than an explicit architectural component. \] 